\chapter{Introduction}

The Deep Kalman Filter (DKF), introduced by Krishnan, Shalit and Sontag
in 2015 \cite{krishnan2015deep}, is a sequential latent-variable model.
It keeps the structure of the Kalman filter, a latent state that evolves
as a Markov chain and emits one observation per time step, but replaces
the linear-Gaussian transition and emission models by Gaussian
distributions whose parameters are computed by neural networks. Since
exact posterior inference is then intractable, the model is trained like
a Variational Auto-Encoder (VAE): an inference network approximates the
posterior over the latent trajectory, and the parameters of both networks
are learned jointly by maximising the Evidence Lower Bound (ELBO).\\

This report presents the mathematics of the DKF, my implementation of
it\footnote{The code is available
\href{https://github.com/CL-BZH/Deep-Kalman-Filters/tree/main}{on
GitHub}.}, and experiments on noisy sine waves and on sequences generated
by the Lorenz attractor\footnote{The cover page shows a realisation of
the Lorenz attractor used to test the models.}. It is organised as
follows.
\begin{itemize}
    \item Chapter 2 recalls the Kalman filter. It is the linear-Gaussian
    special case in which the posterior that the DKF approximates can be
    computed exactly.
    \item Chapter 3 recalls the VAE and, in particular, \emph{why} the
    ELBO is a lower bound on the log-likelihood.
    \item Chapter 4 is the core of the report: the DKF model, the
    structure of its true posterior (derived with d-separation), the
    factorised ELBO, how that ELBO is estimated by sampling, the
    implementation, and three extensions I designed: the autoregressive
    DKF (AR DKF) and its two variants FAR DKF and BAR DKF.
    \item The appendices give the complete derivations and the
    experimental results. The annex recalls d-separation.
\end{itemize}
